\documentclass[conference]{IEEEtran}

\IEEEoverridecommandlockouts

\usepackage{cite}
\usepackage{amsmath,amssymb,amsfonts}
\usepackage{algorithmic}
\usepackage{graphicx}
\usepackage{textcomp}
\usepackage{xcolor}
\usepackage{booktabs}
\usepackage{url}

\def\BibTeX{{\rm B\kern-.05em{\sc i\kern-.025em b}\kern-.08em
    T\kern-.1667em\lower.7ex\hbox{E}\kern-.125emX}}

\begin{document}

\title{GeoRisk-IQ: Geospatial Machine Learning for Bank Branch Opportunity Scoring and Failure-Risk Validation}

\author{
\IEEEauthorblockN{Author Name}
\IEEEauthorblockA{
\textit{Department or Organization}\\
City, Country\\
email@example.com}
}

\maketitle

\begin{abstract}
Financial institutions need interpretable geospatial intelligence for branch
placement, market coverage, and risk monitoring. This paper presents
GeoRisk-IQ, a spatial machine-learning framework that converts point-level
banking data into a grid-based opportunity and risk surface. The system
combines Federal Deposit Insurance Corporation (FDIC) institution records,
financial indicators, and bank-failure labels with local and neighborhood
density features derived from an H3-style spatial grid. GeoRisk-IQ supports
two related tasks: (i) ranking underserved geographic cells as branch
opportunity zones and (ii) validating spatial-financial signals against
historical bank-failure outcomes. In experiments on 27,832 FDIC institution
records aggregated into 6,086 spatial cells, the opportunity scorer achieved
a mean absolute error of 0.0081 and an $R^2$ of 0.9991 against a spatial
proxy target. For failure-risk validation, the hybrid spatial-financial
gradient boosting model achieved a 10-fold cross-validation ROC-AUC of
0.8341 with a bootstrap 95\% confidence interval of [0.8226, 0.8454].
Results suggest that neighborhood spatial structure is a useful signal for
location intelligence and risk validation in banking.
\end{abstract}

\begin{IEEEkeywords}
geospatial machine learning, banking analytics, FDIC, branch location,
spatial risk, opportunity scoring, gradient boosting
\end{IEEEkeywords}

\section{Introduction}
Branch networks remain an important component of retail and community
banking, even as digital channels expand. Decisions about where to expand,
consolidate, or monitor branch infrastructure require more than simple counts
of existing offices. A location can be underserved when it has low local
supply but is near high-demand corridors, dense financial activity, or
adjacent service gaps. Similarly, regional risk may be partially reflected in
spatial patterns that interact with bank-level financial indicators.

GeoRisk-IQ addresses this problem through a reproducible geospatial machine
learning pipeline. The system ingests FDIC institution and financial data,
maps records to a spatial grid, engineers local and neighboring density
features, and trains models for opportunity scoring and failure-risk
validation. The same architecture is designed to generalize to other shortage
or access domains, such as health-professional shortage areas, although this
paper focuses on the banking-domain implementation and validation.

The main contributions are:
\begin{itemize}
    \item A grid-based feature engineering pipeline for banking location
    intelligence using FDIC institution coordinates and financial attributes.
    \item A branch opportunity scoring method that rewards low local supply
    near high surrounding density.
    \item An IEEE-style validation design using historical bank-failure
    labels, cross-validation, temporal holdout testing, bootstrap confidence
    intervals, and feature importance analysis.
    \item An interactive Streamlit implementation for exploring opportunity
    zones, density maps, and validation outputs.
\end{itemize}

\section{Related Work}
Spatial indexing and aggregation are common in transportation, urban
analytics, epidemiology, and financial services. Hexagonal indexing systems
such as H3 are widely used because they support locality-preserving
aggregation and neighborhood operations. GeoRisk-IQ follows this general
principle through an H3-style grid abstraction, allowing point observations
to be converted into cell-level features.

Machine learning has also been used for credit risk, bank distress
prediction, and market opportunity analysis. However, many financial-risk
models focus primarily on institution-level balance-sheet signals. In
contrast, GeoRisk-IQ explicitly represents spatial context through local
cell density, neighboring-cell density, and spatial ratios. This makes the
pipeline suitable for both operational branch planning and empirical
validation against observed bank-failure events.

\section{Data}
\subsection{FDIC Institution Data}
The branch opportunity module uses FDIC institution records containing bank
names, certificates, state abbreviations, cities, coordinates, assets, and
office counts. In the current experiment, 27,832 institutions were loaded
from the FDIC source file \texttt{Geospatiallocation.txt}. Records without
valid latitude or longitude were removed.

\subsection{FDIC Failure Data}
The validation module uses the FDIC failure file
\texttt{FDICFailure.txt}. Failure records are linked to institutions through
certificate identifiers when possible, then mapped to spatial cells. The
current validation grid contains 6,086 cells, of which 1,403 contain at
least one historical failure event.

\subsection{FDIC Financial Data}
Financial indicators are loaded from \texttt{FDICFinancialsFull.txt}. The
current implementation uses the latest available quarterly financial
snapshot and aggregates return on assets, net income, and assets to the
spatial-cell level.

\section{Methodology}
\subsection{Spatial Grid Construction}
Each institution coordinate $(lat_i, lon_i)$ is assigned to a deterministic
spatial cell. Let $c_i$ denote the resulting cell identifier. For each cell
$c$, the branch density is
\begin{equation}
    B_c = \sum_i \mathbb{I}(c_i = c),
\end{equation}
where $B_c$ is the count of institutions in cell $c$. If asset values are
available, total assets are aggregated as
\begin{equation}
    A_c = \sum_i A_i \mathbb{I}(c_i = c).
\end{equation}

\subsection{Neighborhood Features}
For each grid cell, GeoRisk-IQ computes neighboring density averages over
multiple rings. For a neighborhood radius $k$, let $\mathcal{N}_k(c)$ be the
set of neighboring cells within $k$ grid steps of cell $c$. The average
neighbor density is
\begin{equation}
    \bar{B}_{c,k} =
    \frac{1}{|\mathcal{N}_k(c)|}
    \sum_{j \in \mathcal{N}_k(c)} B_j.
\end{equation}
The implemented feature set includes $k \in \{1,2,3\}$, density ratios, log
branch counts, log asset totals, and an isolation indicator.

\subsection{Opportunity Target}
The branch opportunity task uses a proxy target that identifies cells with
low current supply and high surrounding density:
\begin{equation}
    O_c = \frac{\bar{B}_{c,2}}{1 + B_c}.
\end{equation}
This formulation rewards locations that are near dense banking corridors but
have relatively low local branch counts. A gradient boosting regressor is
trained on engineered cell features, and predicted scores are normalized to
a 0--100 opportunity scale.

\subsection{Failure-Risk Validation}
For risk validation, each cell receives a binary label:
\begin{equation}
    y_c =
    \begin{cases}
    1, & \text{if at least one bank failure is mapped to } c,\\
    0, & \text{otherwise.}
    \end{cases}
\end{equation}
Four model families are compared:
\begin{itemize}
    \item logistic regression using the full hybrid feature set,
    \item spatial-only gradient boosting,
    \item financial-only gradient boosting,
    \item hybrid spatial-financial gradient boosting.
\end{itemize}
Evaluation uses 10-fold stratified cross-validation, temporal holdout
testing, bootstrap confidence intervals, permutation tests, and feature
importance analysis.

\section{Experimental Setup}
The pipeline was executed in Python using pandas, NumPy, scikit-learn,
Plotly, pydeck, Streamlit, SciPy, and SHAP. Gradient boosting models use a
fixed random seed of 42. The opportunity scorer uses a train-test split with
20\% held out. Failure-risk validation uses 10-fold stratified
cross-validation and a temporal split that trains on pre-2010 failures and
tests on post-2010 failures where sufficient positive examples exist.

\section{Results}
\subsection{Opportunity Scoring}
The opportunity regression task produced strong agreement with the proxy
target, as shown in Table~\ref{tab:opportunity_metrics}. The top-ranked
cells represent areas with low local branch count but elevated surrounding
activity.

\begin{table}[htbp]
\caption{Opportunity Scoring Performance}
\begin{center}
\begin{tabular}{lc}
\toprule
\textbf{Metric} & \textbf{Value}\\
\midrule
Institutions loaded & 27,832\\
Spatial cells & 6,086\\
Mean absolute error & 0.0081\\
$R^2$ & 0.9991\\
\bottomrule
\end{tabular}
\label{tab:opportunity_metrics}
\end{center}
\end{table}

\begin{table}[htbp]
\caption{Example Top Opportunity Cells}
\begin{center}
\begin{tabular}{lccc}
\toprule
\textbf{Cell Center} & \textbf{Score} & \textbf{Branches} & \textbf{Assets}\\
\midrule
(41.4, -88.0) & 100.0 & 1 & 260,568\\
(29.4, -95.4) & 84.4 & 1 & 12,268\\
(32.6, -96.6) & 82.9 & 1 & 376,443\\
(30.0, -95.8) & 75.8 & 1 & 27,025\\
(39.2, -77.0) & 75.4 & 1 & 13,765,535\\
\bottomrule
\end{tabular}
\label{tab:top_opportunities}
\end{center}
\end{table}

\subsection{Failure-Risk Cross-Validation}
Table~\ref{tab:model_comparison} reports cross-validation results. The
hybrid spatial-financial model achieved the highest ROC-AUC, while logistic
regression achieved the highest F1 score in the tested configuration. This
suggests that hybrid features improve ranking quality, while classification
thresholding and class imbalance remain important design considerations.

\begin{table}[htbp]
\caption{10-Fold Cross-Validation Model Comparison}
\begin{center}
\begin{tabular}{lcc}
\toprule
\textbf{Model} & \textbf{ROC-AUC} & \textbf{F1}\\
\midrule
Logistic regression & 0.8143 & 0.5780\\
Spatial-only GBM & 0.8279 & 0.5354\\
Financial-only GBM & 0.6029 & 0.1697\\
Hybrid GBM & 0.8341 & 0.5320\\
\bottomrule
\end{tabular}
\label{tab:model_comparison}
\end{center}
\end{table}

\subsection{Temporal Holdout}
The temporal holdout split trained on pre-2010 failures and clean cells,
then tested on post-2010 failures and held-out clean cells. The hybrid model
achieved the results shown in Table~\ref{tab:temporal}.

\begin{table}[htbp]
\caption{Temporal Holdout Results}
\begin{center}
\begin{tabular}{lc}
\toprule
\textbf{Metric} & \textbf{Value}\\
\midrule
Precision & 0.3776\\
Recall & 0.4512\\
F1 & 0.4111\\
ROC-AUC & 0.7331\\
\bottomrule
\end{tabular}
\label{tab:temporal}
\end{center}
\end{table}

\subsection{Statistical Tests and Feature Importance}
The hybrid model achieved a bootstrap 95\% confidence interval of
[0.8226, 0.8454] for cross-validated ROC-AUC. A permutation test comparing
the hybrid model with the spatial-only model gave $p=0.0013$, suggesting a
statistically meaningful AUC improvement. The strongest SHAP feature in the
current run was the second-ring density ratio, followed by first- and
second-ring neighbor density and mean financial assets.

\section{Discussion}
The results indicate that spatial context provides meaningful information
for banking location intelligence. The opportunity scorer surfaces cells
with low local branch count that are near higher-density corridors, while
the validation pipeline shows that neighborhood structure contributes to
bank-failure ranking performance. The relatively weak financial-only model
in this setup may reflect the use of a latest-quarter aggregate rather than
full historical financial time series. Future versions should compare
quarter-specific, lagged, and rolling financial features.

The high $R^2$ in opportunity scoring should be interpreted carefully because
the target is a deterministic proxy derived from spatial features. Its role
is to rank candidate opportunity zones rather than to infer observed branch
openings. The failure-risk validation experiment provides a stronger
external check because it uses historical event labels.

\section{Limitations}
Several limitations remain. First, the current spatial grid is an H3-style
deterministic abstraction rather than a production H3 library dependency.
Second, bank failures are mapped through institution certificates and
available coordinates; failures without reliable coordinates are omitted.
Third, opportunity scoring uses a proxy target because true expansion-demand
labels are not directly available. Fourth, the current financial feature set
uses aggregate cell-level indicators and should be extended with historical
lags.

\section{Conclusion}
GeoRisk-IQ demonstrates how geospatial machine learning can support bank
branch opportunity scoring and risk validation. By combining local supply,
neighborhood density, and financial indicators, the framework produces
interpretable opportunity zones and empirically testable risk surfaces. The
current results support further development of spatial-financial models for
branch network planning, regulatory analytics, and cross-domain access
intelligence.

\section*{Acknowledgment}
The author thanks the open-data providers whose datasets make reproducible
geospatial financial analysis possible. This section can be updated with
specific funding, institutional, or collaborator acknowledgments.

\begin{thebibliography}{00}
\bibitem{fdic}
Federal Deposit Insurance Corporation, ``BankFind Suite API,'' FDIC.
[Online]. Available: \url{https://banks.data.fdic.gov/}

\bibitem{h3}
Uber Technologies, ``H3: A Hexagonal Hierarchical Geospatial Indexing
System.'' [Online]. Available: \url{https://h3geo.org/}

\bibitem{friedman}
J. H. Friedman, ``Greedy function approximation: A gradient boosting
machine,'' \textit{Annals of Statistics}, vol. 29, no. 5, pp. 1189--1232,
2001.

\bibitem{shap}
S. M. Lundberg and S.-I. Lee, ``A unified approach to interpreting model
predictions,'' in \textit{Proc. Advances in Neural Information Processing
Systems}, 2017, pp. 4765--4774.

\bibitem{sklearn}
F. Pedregosa et al., ``Scikit-learn: Machine learning in Python,''
\textit{Journal of Machine Learning Research}, vol. 12, pp. 2825--2830,
2011.
\end{thebibliography}

\end{document}
